\chapter*{ABSTRACT}
\addcontentsline{toc}{chapter}{ABSTRACT}
This study developed a TinyML Smart Plug as complementary socket-level
protection for progressive fire-inducing electrical faults in 230V
residential circuits. The prototype addressed sustained overload,
degraded socket-contact heating, and load-limited series arc faults,
which may remain hazardous before conventional breaker operation. It
integrated voltage, current, and temperature sensing with an ESP32-S3
controller, relay-based load isolation, local audio-visual alerts,
Firebase logging, and a Progressive Web App for monitoring and basic
relay control. Local firmware handled calibration, thermal estimation,
threshold protection, load-context classification, and Random Forest
series-arc inference; the cloud interface supported visualization,
logging, labeling, and operator interaction. Evaluation included sensor
calibration and revalidation, controlled overload tests, socket-heating
characterization and protection tests, induced series-arc data
collection, waveform-feature extraction, manual labeling, and embedded
model deployment. Revalidation produced mean absolute errors of 3.910V,
0.178A, and 1.02°C for the voltage, current, and temperature channels,
respectively. The highest-stress overload case reached 17.30A mean
current and 19.73A peak current before relay disconnection after
61.382s. Heating tests showed faster temperature rise in the degraded
socket than in the normal socket under comparable loading. For
trial-level series arc detection, the embedded model and firmware
persistence logic achieved 97.50\% accuracy, 98.97\% precision, 96.00\%
recall, and 97.46\% F1-score across tested loads. The results support
controlled laboratory feasibility for calibrated sensing, deterministic
threshold protection, and on-device TinyML inference in a smart plug that
complements, but does not replace, mandatory residential electrical
protection.

\textbf{Keywords:} TinyML, Smart Plug, Series Arc Fault Detection,
Socket Heating, Overload Protection

